\subsection*{Nilpotent truncation of the cotangent module}

Notation in this file is local to the nilpotent truncation lane.  Taylor
windows in the main formal-Darboux theorem are
\(\mathfrak h^{\mathrm{Taylor}}_{\leq W}=\mathfrak m/\mathfrak m^{W+1}\).
They are finite coefficient projections, not Hamiltonian Lie quotients:
\(\{z_1,z_2^{W+1}\}=(W+1)z_2^W\) shows that
\(\mathfrak m^{W+1}\) is not a Poisson ideal of \(\mathfrak m\).
The cubic nilpotent truncation is
\[
  \mathfrak h^{\mathrm{nilp}}_{[3,N]}
  :=\mathfrak m^3/\mathfrak m^{N+1}.
\]
The Taylor-window symbol is reserved for coefficient projections of the
full \(\mathfrak m\)-adic Hamiltonian source in the formal-Darboux theorem
and its Matlis/Tate completion appendices.

The full CE/PV trace-sector map of Theorem~\ref{thm:main-local} is a
completed formal-Darboux statement for the reduced Hamiltonian Lie algebra
\(\mathfrak h=\mathfrak m=(z_1,z_2)\C[[z_1,z_2]]\).  The lower-central
pronilpotent route is a separate finite nilpotent-window route.  It
starts with the cubic Lie subalgebra
\(\mathfrak m^3\subset\mathfrak h\).  It does not prove the full
formal-Darboux CE/PV map, and it does not obstruct that map.  It records
the obstruction to recovering the full Hamiltonian source from
nilpotent \(\mathfrak m\)-adic quotients.  The obstruction to this
nilpotent recovery is the explicit triple
\(\bigl(\mathfrak o^{\mathrm{lin}}_N,\mathfrak o^{\mathrm{ss}},
\mathfrak o^{\mathrm{Cas}}_{\leq2}\bigr)\).

For every $N\geq 1$,
\[
  \{z_1,z_2^{N+1}\}=(N+1)z_2^N\notin\mathfrak m^{N+1},\qquad
  \{z_1^2,z_2^2\}=4z_1z_2,\qquad [H_2,H_2]=H_2\cong\mathfrak{sl}_2.
\]
The first identity is the linear-tail obstruction
$\mathfrak o^{\mathrm{lin}}_N(z_1;z_2^{N+1})=(N+1)z_2^N$, ruling out
$\mathfrak m^{N+1}$ as a Lie ideal once $z_1\in H_1$ is admitted.  The
last two encode the semisimple obstruction
$\mathfrak o^{\mathrm{ss}}\colon\Lambda^2H_2\to H_2$ with image
$\mathfrak{sl}_2$, ruling out nilpotence of $H_2$.  Together with the
low-degree Casimir component
$\mathfrak o^{\mathrm{Cas}}_{\leq2}=\sum_ie_{1,i}\otimes
e_1^i+\sum_ie_{2,i}\otimes e_2^i$, these form the obstruction triple
of Definition~\ref{defn:nilpotent-low-degree-obstruction-terms}.

On the cubic ideal
$\mathfrak h^{\mathrm{nilp}}_{[3,N]}=\mathfrak m^3/\mathfrak m^{N+1}$, where the
obstruction triple is silent, the finite cotangent CE/PV identity
$\Phi_{\leq N}\colon
\operatorname{Regr}_u
C^\bullet_{\mathrm{CE},\mathrm{cont}}(\mathfrak g^{\mathrm{nilp}}_{[3,N],I})
\xrightarrow{\sim}
\bigl(\mathrm{PV}(\widehat{\mathbf S}(\mathfrak h^{\mathrm{nilp}}_{[3,N],I})),
d_{\pi_N}\bigr)$ is the
cubic-truncation trace-sector model for
\(\mathcal A^{\mathrm{cl}}_{\mathrm{bulk}}|_b\), proved in the finite
classical sector
(Theorem~\ref{thm:phi-trunc-classical}).
Here \(\pi_N\) is the linear Kirillov Poisson bivector determined by
the finite Lie bracket on \(\mathfrak h^{\mathrm{nilp}}_{[3,N]}\).
Theorem~\ref{thm:phi-hbar-trunc} excludes the formal Moyal lift on
the same finite truncation;
Theorem~\ref{thm:colim-obstruction-full-h} shows that the inverse limit
of this nilpotent tower recovers the proper subalgebra
\(\mathfrak m^3\subsetneq\mathfrak h\).  The full formal-Darboux CE/PV
map is given elsewhere by the completed coordinate construction.
Problem~\ref{prob:linear-sector-lift} is the narrower problem of
replacing the nilpotent finite-window route by a relative or reductive
route that includes \(H_1\oplus H_2\) without using false
pronilpotence.

\subsection*{The truncation}

Let $\mathfrak h=(z_1,z_2)\C[[z_1,z_2]]=\mathfrak m$ be the reduced
Hamiltonian Lie algebra of the formal symplectic polydisk on $\C^2$, with
bracket $\{f,g\}=\partial_{z_1}f\,\partial_{z_2}g-
\partial_{z_2}f\,\partial_{z_1}g$ modulo constants. Let
$\mathfrak m^k=(z_1,z_2)^k\C[[z_1,z_2]]$ for $k\geq 1$ be the
$k$-th power of the maximal ideal. Fix integers $N\geq k\geq 3$ and
define the \emph{nilpotent truncation}
$$\mathfrak h^{\mathrm{nilp}}_{[k,N]}:=\mathfrak m^k/\mathfrak m^{N+1},$$
with $\mathfrak m$-adic Poisson bracket inherited from the
ambient $\mathfrak h_{\mathrm{poly}}$. The truncation used below is
$k=3$:
$$\mathfrak h^{\mathrm{nilp}}_{[3,N]}=
  \mathfrak m^3/\mathfrak m^{N+1}.$$
Concretely, $\mathfrak h^{\mathrm{nilp}}_{[3,N]}$ is the finite-dimensional vector
space spanned by monomials $z_1^az_2^b$ with $3\leq a+b\leq N$.

The cutoff lies at $k=3$.  At $k=2$, the degree-$2$ piece
$H_2=\mathrm{span}\{z_1^2,z_1z_2,z_2^2\}$ closes under the Poisson
bracket and carries the $\mathfrak{sl}_2$-relations
$$\{z_1^2,z_2^2\}=4z_1z_2,\quad
  \{z_1^2,z_1z_2\}=2z_1^2,\quad
  \{z_1z_2,z_2^2\}=2z_2^2,$$
so $\mathfrak m^2/\mathfrak m^{N+1}$ contains a nonabelian semisimple
subalgebra and violates nilpotence, hence the lower-central
Tate-pronilpotence hypothesis of \ref{lem:continuous-bar-cobar}.  At
$k=3$ the bracket raises $\mathfrak m$-adic degree strictly.

\begin{lem}[Truncation is a finite-dimensional nilpotent Lie algebra]
\label{lem:trunc-is-nilpotent-lie}
For every $N\geq 3$:
\begin{enumerate}
  \item $\mathfrak h^{\mathrm{nilp}}_{[3,N]}=\mathfrak m^3/\mathfrak m^{N+1}$ is a
    finite-dimensional Lie algebra of dimension
    $\binom{N+2}{2}-6$ (the number of
    monomials of degree between $3$ and $N$ inclusive).
  \item $\mathfrak h^{\mathrm{nilp}}_{[3,N]}$ is nilpotent of class at most
    $N-2$ in the lower central series.
  \item $(\mathfrak h^{\mathrm{nilp}}_{[3,N]})^\vee$ is finite-dimensional, hence the
    formal Casimir
    $C_{\mathfrak h^{\mathrm{nilp}}_{[3,N]}}=\sum_{3\leq d\leq N}\sum_i
     e_{d,i}\otimes e_d^i$
    is a convergent element of
    $\mathfrak h^{\mathrm{nilp}}_{[3,N]}\otimes(\mathfrak h^{\mathrm{nilp}}_{[3,N]})^\vee$ as a finite
    sum in the ordinary tensor product.
  \item The linear Hamiltonians $z_1,z_2\in\mathfrak m\setminus
    \mathfrak m^3$, the conformal-symplectic generators
    $H_2=\mathrm{span}\{z_1^2,z_1z_2,z_2^2\}\cong\mathfrak{sl}_2$,
    and the constant Hamiltonian $1\notin\mathfrak m$ are omitted
    from $\mathfrak h^{\mathrm{nilp}}_{[3,N]}$ by construction.
\end{enumerate}
\end{lem}

\begin{proof}
\emph{(i)} The space $\mathfrak m^3$ is a Lie subalgebra of the
polynomial Hamiltonian Lie algebra $\mathfrak h_{\mathrm{poly}}$, and
$\mathfrak m^{N+1}$ is a Lie ideal inside $\mathfrak m^3$: for
$f\in\mathfrak m^p$ and
$g\in\mathfrak m^q$ one has $\{f,g\}\in\mathfrak m^{p+q-2}$, since the
Poisson bracket on $\C[z_1,z_2]$ lowers total polynomial degree by
exactly two, and each partial derivative lowers $\mathfrak m$-adic
degree by exactly one. For $p,q\geq 3$ the bracket lands in
$\mathfrak m^{p+q-2}\subset\mathfrak m^4\subset\mathfrak m^3$, making
$\mathfrak m^3$ a Lie subalgebra; for $p\geq 3$ and $q\geq N+1$
the bracket lands in $\mathfrak m^{p+q-2}\subset\mathfrak m^{N+2}
\subset\mathfrak m^{N+1}$, making $\mathfrak m^{N+1}$ a Lie ideal of
$\mathfrak m^3$. The quotient
$\mathfrak h^{\mathrm{nilp}}_{[3,N]}=\mathfrak m^3/\mathfrak m^{N+1}$ inherits a
well-defined Lie bracket. Dimension count: monomials $z_1^az_2^b$
with $3\leq a+b\leq N$ number $\binom{N+2}{2}-6$ (the total number
with $0\leq a+b\leq N$ minus the six of degrees $0,1,2$).

\emph{(ii)} The Poisson bracket of two elements of
$\mathfrak m^3$ raises $\mathfrak m$-adic degree by at least $1$:
$\{H_p,H_q\}\subset H_{p+q-2}\subset\mathfrak m^{p+q-2}$, and for
$p,q\geq 3$, $p+q-2\geq 4>p$ (since $q\geq 3$). Thus the descending
Lie series
$$\bigl(\mathfrak h^{\mathrm{nilp}}_{[3,N]}\bigr)^{(\geq m)}:=
  \underbrace{[\mathfrak h^{\mathrm{nilp}}_{[3,N]},[\mathfrak h^{\mathrm{nilp}}_{[3,N]},\ldots]]}_{m\text{ brackets}}
  \subset\bigl(\mathfrak m^{m+3}/\mathfrak m^{N+1}\bigr)$$
satisfies $\bigl(\mathfrak h^{\mathrm{nilp}}_{[3,N]}\bigr)^{(\geq m)}=0$ for
$m+3\geq N+1$, i.e.\ for $m\geq N-2$. Hence $\mathfrak h^{\mathrm{nilp}}_{[3,N]}$
is nilpotent of class at most $N-2$.

\emph{(iii)} Immediate from finite-dimensionality.

\emph{(iv)} Immediate from the $\mathfrak m^3$-adic cutoff and the
exclusion $a+b\leq 2$.
\end{proof}

The truncation kills the linear and conformal-symplectic sectors:
the obstruction
$\{z_1,z_2^{N+1}\}=(N+1)z_2^N\notin\mathfrak m^{N+1}$ of
\ref{lem:finite-window-hamiltonian-obstruction} requires $z_1\in H_1$,
which the cubic cutoff omits.  The Casimir is the finite sum
$C_{\mathfrak h^{\mathrm{nilp}}_{[3,N]}}=\sum_{d=3}^N\sum_i e_{d,i}\otimes e_d^i\in
\mathfrak h^{\mathrm{nilp}}_{[3,N]}\otimes(\mathfrak h^{\mathrm{nilp}}_{[3,N]})^\vee$; the
product/direct-sum support obstruction of
\ref{lem:tate-casimir-obstruction}, which depends on infinite degree
expansion, vanishes.

\subsection*{The classical truncated cotangent CE/PV identification}

\begin{thm}[Classical truncated CE/PV identification: Koszul resolution on the cubic ideal]\label{thm:phi-trunc-classical}
Set
$\mathfrak g^{\mathrm{nilp}}_{[3,N]}=\mathfrak h^{\mathrm{nilp}}_{[3,N]}\ltimes
(\mathfrak h^{\mathrm{nilp}}_{[3,N]})^\vee[1]$. For every interval $I\subset\R$ let
$\mathfrak g^{\mathrm{nilp}}_{[3,N],I}=
\Omega^\bullet_c(I)\,\widehat\otimes\,\mathfrak h^{\mathrm{nilp}}_{[3,N]}\,\ltimes\,
\mathcal D'_c(I)\,\widehat\otimes\,(\mathfrak h^{\mathrm{nilp}}_{[3,N]})^\vee[1]$.
Then, after the Poisson, residue, and interval-smearing conventions fixed
above, the coordinate dictionary is
\[
  c^\rho\longmapsto\partial_\rho,
  \qquad
  u_f\longmapsto O_f,
  \qquad
  d_{\mathrm{CE}}u_f\longmapsto X_f=d_{\pi_N}O_f.
\]
Here the standard Chevalley--Eilenberg degree of \(u_f\) is \(2\),
whereas its coordinate degree after \(\operatorname{Regr}_u\) is \(0\);
\(|c^\rho|=|\partial_\rho|=1\).  This is the coordinate Koszul model of
\(\mathcal A^{\mathrm{cl}}_{\mathrm{bulk}}|_b\) on the cubic
truncation: a filtered cochain isomorphism of completed dg
graded-commutative \(P_0\)-algebras
$$\Phi_{\leq N}:
  \operatorname{Regr}_u
  C^\bullet_{\mathrm{CE},\mathrm{cont}}(\mathfrak g^{\mathrm{nilp}}_{[3,N],I})
  \xrightarrow{\ \sim\ }
  \bigl(
  \mathrm{PV}\bigl(\widehat{\mathbf S}(\mathfrak h^{\mathrm{nilp}}_{[3,N],I})\bigr),
  d_{\pi_N}\bigr)$$
on every interval \(I\subset\R\), with the same Schouten/Lichnerowicz
convention as in \ref{thm:reduced-ce-pv-central-operation}. It is
compatible with extension by zero and with disjoint-interval products in
this reduced finite model, and on the Hamiltonian generators it recovers
the stalkwise operation \(f\mapsto L_{B_f}\). The cotangent generators
are finite principal-part labels; their unreduced boundary centrality
homotopies remain part of
Problem~\ref{prob:formal-factorization-center}.  An identification with
ordinary or chiral Hochschild cochains is not part of this finite
CE/PV identity: ordinary Hochschild--Kostant--Rosenberg comparison has
zero polyvector differential, while the displayed target has
\(d_{\pi_N}\).  A Hochschild target therefore requires the supplied
filtered twisted-formality and Hochschild-centre lift data.
\end{thm}

\begin{proof}
The Lie algebra $\mathfrak g^{\mathrm{nilp}}_{[3,N]}$ is finite-dimensional
($\dim\mathfrak g^{\mathrm{nilp}}_{[3,N]}=2\dim\mathfrak h^{\mathrm{nilp}}_{[3,N]}=
2\binom{N+2}{2}-12$), hence pro-finite-dimensional in the trivial
constant filtration. It is lower-central Tate-pronilpotent in the sense of
\ref{lem:continuous-bar-cobar}: by
\ref{lem:trunc-is-nilpotent-lie}(ii) the descending Lie powers of
$\mathfrak h^{\mathrm{nilp}}_{[3,N]}$ vanish after at most $N-2$ brackets.  The
cotangent summand is abelian, and nilpotence of the semidirect product
uses the degree-lowering
coadjoint action: since \([H_p,H_q]\subset H_{p+q-2}\), an element
\(f\in H_p\), \(p\geq3\), sends \(H_d^\vee\) to
\(H_{d-p+2}^\vee\), hence lowers dual degree by at least one.  Repeated
action therefore kills every finite quotient
\((\mathfrak h^{\mathrm{nilp}}_{[3,N]})^\vee\).  Together with nilpotence of
\(\mathfrak h^{\mathrm{nilp}}_{[3,N]}\), the descending Lie powers of
\(\mathfrak g^{\mathrm{nilp}}_{[3,N]}=\mathfrak h^{\mathrm{nilp}}_{[3,N]}\ltimes
(\mathfrak h^{\mathrm{nilp}}_{[3,N]})^\vee[1]\) also terminate. Hence
$\bigcap_m\bigl(\mathfrak g^{\mathrm{nilp}}_{[3,N]}\bigr)^{(\geq m)}=0$ in any Tate
filtration, and the lower-central Tate-pronilpotence hypothesis of
\ref{lem:continuous-bar-cobar} is satisfied trivially.

Apply \ref{thm:reduced-ce-pv-central-operation} (the cochain-level reduced
Hamiltonian CE/PV central-operation theorem) to $\mathfrak h^{\mathrm{nilp}}_{[3,N]}$.
The conclusion is the completed dg graded-commutative \(P_0\)-algebra
isomorphism specified by that theorem,
$$\Phi_{\leq N}^{\mathrm{stalk}}:
  \operatorname{Regr}_u
  C^\bullet_{\mathrm{CE},\mathrm{cont}}(\mathfrak g^{\mathrm{nilp}}_{[3,N]})
  \xrightarrow{\ \sim\ }
  \bigl(
  \mathrm{PV}\bigl(\widehat{\mathbf S}(\mathfrak h^{\mathrm{nilp}}_{[3,N]})\bigr),
  d_{\pi_N}\bigr),$$
with the explicit generator-level data of
\ref{thm:reduced-ce-pv-central-operation}.

To pass from the stalk identification to the factorization-algebra
identification on intervals $I\subset\R$, the Hamiltonian-sector
factorization extension of \ref{thm:hamiltonian-current-factorization}
applies verbatim with $\mathfrak h$ replaced by
$\mathfrak h^{\mathrm{nilp}}_{[3,N]}$: the locality bracket
$\{(\phi_1)^i{}_j(t),(\phi_2)^k{}_l(t')\}=\delta^i_l\delta^k_j\delta(t-t')$
on two brane-time points is unchanged, the smearing
  $\mathfrak h^{\mathrm{nilp}}_{[3,N],I}=\Omega^\bullet_c(I)\,\widehat\otimes\,
  \mathfrak h^{\mathrm{nilp}}_{[3,N]}$ uses the ordinary tensor product because
\(\mathfrak h^{\mathrm{nilp}}_{[3,N]}\) is finite-dimensional, and the strict \(P_0\)
bracket follows from the delta-function locality.

In this finite truncated classical sector, the cotangent factor has a finite
Tate-Casimir kernel. The principal-part sub-factorization algebra
$$\widetilde{\mathrm{Obs}}^{\mathrm{cl,pp}}_{\partial,\leq N}(I)
  =\mathrm{Obs}^{\mathrm{cl}}_{\mathrm{open},\leq N}(I)
   \,\widehat\otimes\,
   \widehat{\mathbf S}\bigl((\mathfrak h^{\mathrm{nilp}}_{[3,N],I})^\vee[1]\bigr)$$
    admits a chain-level primitive projection
    $\pi_{\mathrm{prim}}^{\leq N}$ because the Tate Casimir
    $C_{\mathfrak h^{\mathrm{nilp}}_{[3,N]}}=\sum_{d=3}^N\sum_i e_{d,i}\otimes e_d^i$
is a finite sum, hence a literal element of the (uncompleted) tensor
product $\mathfrak h^{\mathrm{nilp}}_{[3,N]}\otimes(\mathfrak h^{\mathrm{nilp}}_{[3,N]})^\vee$. The
$Q$-compatibility, factorization compatibility, and $P_0$-bracket
homotopies are then formal, since every step uses finite-support tensors
in the direct-sum dual.

The bracket compatibility and graded Leibniz extension to the completed
finite CE/PV algebra is the same argument as in
\ref{thm:reduced-ce-pv-central-operation}, with the additional
simplification that all completions are trivial because
$\mathfrak h^{\mathrm{nilp}}_{[3,N]}$ is finite-dimensional. Compatibility with
extension by zero, disjoint-interval factorization, and the
stalkwise limit follow as in
\ref{thm:hamiltonian-current-factorization}.
\end{proof}

\subsection*{Quantum obstruction for the finite truncation}

\begin{defn}[Moyal-flat nilpotent replacement datum]
\label{defn:nilpotent-moyal-flat-datum}
Let \(N\geq4\) and put
\[
  V_N=\mathfrak h^{\mathrm{nilp}}_{[3,N]}
  =\mathfrak m^3/\mathfrak m^{N+1},\qquad
  Q_{<3}=\mathfrak m/\mathfrak m^3 .
\]
Let \(\pi_{[3,N]}\) denote projection to monomials of total degree
between \(3\) and \(N\), and let \(\pi_{<3}\) denote projection to
\(Q_{<3}\).  For \(f,g\in V_N\), choose representatives
\(\tilde f,\tilde g\in\mathfrak m^3\).  The Moyal closure cochain is
\[
  \operatorname{cl}_{N,\hbar}(f,g)
  =
  \pi_{<3}\bigl(\{\tilde f,\tilde g\}_{\hbar_{\mathrm{W}}}\bigr)
  \in Q_{<3}[[\hbar_{\mathrm{W}}]],
\]
an element of
\(\operatorname{Hom}^2(\Lambda^2V_N,Q_{<3})[[\hbar_{\mathrm{W}}]]\).
The projected Moyal bracket on \(V_N[[\hbar_{\mathrm{W}}]]\) is
\[
  [f,g]^{\operatorname{proj}}_{N,\hbar}
  =
  \pi_{[3,N]}\bigl(\{\tilde f,\tilde g\}_{\hbar_{\mathrm{W}}}\bigr),
\]
and its Jacobi obstruction is
\[
  \operatorname{Jac}_{N,\hbar}(f,g,k)
  =
  [f,[g,k]^{\operatorname{proj}}_{N,\hbar}]^{\operatorname{proj}}_{N,\hbar}
  +\textup{cyclic}
  \in V_N[[\hbar_{\mathrm{W}}]] .
\]
The cochain-compatibility obstruction
\(\operatorname{CE}_{N,\hbar}\) is the degree-two endomorphism
\[
  (d^{\operatorname{proj}}_{\mathrm{CE},N,\hbar})^2
  \in
  \operatorname{End}^2
  C^\bullet_{\mathrm{CE}}\bigl(
    V_N[[\hbar_{\mathrm{W}}]]
    \ltimes V_N^\vee[[\hbar_{\mathrm{W}}]][1]\bigr),
\]
where \(d^{\operatorname{proj}}_{\mathrm{CE},N,\hbar}\) is built from
the projected bracket and the induced coadjoint action.

A Moyal-flat nilpotent replacement consists of a
\(\C[[\hbar_{\mathrm{W}}]]\)-module \(V^{\flat}_{N,\hbar}\), maps
\[
  V_N[[\hbar_{\mathrm{W}}]]
  \xrightarrow{\ i_N\ }
  V^{\flat}_{N,\hbar}
  \xrightarrow{\ q_N\ }
  V_N[[\hbar_{\mathrm{W}}]],
  \qquad q_Ni_N=\operatorname{id},
\]
and a Lie bracket on \(V^{\flat}_{N,\hbar}\) whose projection to
\(V_N\) has vanishing
\(\operatorname{cl}_{N,\hbar}\), \(\operatorname{Jac}_{N,\hbar}\), and
\(\operatorname{CE}_{N,\hbar}\).  Equivalently, a strict projected
finite quantum source on \(V_N\) is allowed only after these three
obstructions vanish for the projected bracket and its coadjoint action.
\end{defn}

\begin{thm}[Quantum obstruction for the nilpotent truncation]
\label{thm:phi-hbar-trunc}
  The classical finite truncation
  \ref{thm:phi-trunc-classical} is proved in the finite classical
  sector.  The same finite
  Lie algebra
  $\mathfrak h^{\mathrm{nilp}}_{[3,N]}=\mathfrak m^3/\mathfrak m^{N+1}$ has an
  obstruction to a formal Moyal deformation for $N\geq4$: the
  Moyal closure cochain of
  Definition~\ref{defn:nilpotent-moyal-flat-datum} is nonzero.
  Consequently a formal quantum descendant theorem on the strict finite
  truncation requires the Moyal-flat nilpotent replacement datum of
  Definition~\ref{defn:nilpotent-moyal-flat-datum}, or a projected
  bracket for which
  \(\operatorname{Jac}_{N,\hbar}=0\) and
  \(\operatorname{CE}_{N,\hbar}=0\).

  What remains true at quantum level is the restriction of the
  full-sector degree-zero formal Moyal map
  \ref{thm:phi-hbar-all-order} to the finite set of generators of
  $\mathfrak h^{\mathrm{nilp}}_{[3,N]}$, viewed inside the full
  $\mathfrak h_{\hbar_{\mathrm{W}}}$.  This is a restriction of a proved map on the
  full Hamiltonian source. A closed CE theorem for a finite quantum Lie
  subalgebra is equivalent to the vanishing of the closure, Jacobi, and
  CE-square obstructions in
  Definition~\ref{defn:nilpotent-moyal-flat-datum}.
\end{thm}

\begin{proof}
  The Moyal bracket is
  $$\{f,g\}_{\hbar_{\mathrm{W}}}=\sum_{j\geq 0}
    \frac{\hbar_{\mathrm{W}}^{2j}}{(2j+1)!\,2^{2j}}P^{2j+1}(f,g),\qquad
    P=\partial_{z_1}\otimes\partial_{z_2}
      -\partial_{z_2}\otimes\partial_{z_1}.$$
  If $f\in\mathfrak m^p$ and $g\in\mathfrak m^q$, the
  order-\(\hbar_{\mathrm{W}}^{2j}\) summand lies in
  $\mathfrak m^{p+q-2(2j+1)}$ when nonzero.  The Poisson term
  $j=0$ preserves $\mathfrak m^3$ as a Lie subalgebra, but the
  higher Moyal terms can leave the ideal.  For $N\geq4$,
  $$P^3(z_1^4,z_2^4)=24z_1\cdot24z_2=576z_1z_2,$$
  so
  $$\{z_1^4,z_2^4\}_{\hbar_{\mathrm{W}}}
    =16z_1^3z_2^3+24\hbar_{\mathrm{W}}^2 z_1z_2+\cdots.$$
  The term \(24\hbar_{\mathrm{W}}^2z_1z_2\) has \(\mathfrak m\)-adic degree \(2\).
  It has degree \(2\) and lies outside both the cubic ideal
  \(\mathfrak m^3\) and the constants killed by the reduction modulo
  \(\C\cdot1\).  Therefore
  \[
    \operatorname{cl}_{N,\hbar}(z_1^4,z_2^4)
    =
    24\hbar_{\mathrm{W}}^2z_1z_2
    \neq0
    \quad\text{in }Q_{<3}[[\hbar_{\mathrm{W}}]],
  \]
  and \(\mathfrak m^3/\mathfrak m^{N+1}\) has a Moyal-closure
  obstruction for \(N\geq4\).

  Therefore the CE differential for a hypothetical finite quantum
  source
  $\mathfrak h^{\mathrm{nilp}}_{[3,N],\hbar_{\mathrm{W}}}\ltimes
  (\mathfrak h^{\mathrm{nilp}}_{[3,N],\hbar_{\mathrm{W}}})^\vee[1]$ is incompatible with simple
  restriction of the Moyal bracket.  A projected bracket carries the
  Jacobi class \(\operatorname{Jac}_{N,\hbar}\) and the CE-square class
  \(\operatorname{CE}_{N,\hbar}\); the projected-degree warning of
  \ref{lem:finite-window-hamiltonian-obstruction} shows that this is
  a genuine Jacobi condition.  The formal Moyal theorem
  \ref{thm:phi-hbar-all-order} remains valid on the full
  Hamiltonian algebra $\mathfrak h_{\hbar_{\mathrm{W}}}$, and its map can be
  evaluated on generators of degrees $3,\ldots,N$.  That operation yields
  restricted values; a finite quantum CE source requires the vanishing
  data of Definition~\ref{defn:nilpotent-moyal-flat-datum}.  The descendant
  cotangent sector at \(\hbar_{\mathrm{W}}=0\) is exactly
  \ref{thm:phi-trunc-classical}; at \(\hbar_{\mathrm{W}}\neq0\) it is the named
  Moyal-flat truncation problem.
\end{proof}

\begin{prop}[Sectors removed by truncation]\label{prop:trunc-lost-sectors}
The truncation $\mathfrak h\rightsquigarrow\mathfrak h^{\mathrm{nilp}}_{[3,N]}=
\mathfrak m^3/\mathfrak m^{N+1}$ removes the following five
data from the local theorem \ref{thm:main-local}:
\begin{enumerate}
  \item\label{lost:translations}
    The linear Hamiltonians $z_1,z_2\in\mathfrak m\setminus\mathfrak m^2$,
    which generate translations in the symplectic plane and the
    center-of-mass mode of the brane.
  \item\label{lost:sl2}
    The conformal-symplectic $\mathfrak{sl}_2$ generators
    $H_2=\mathrm{span}\{z_1^2,z_1z_2,z_2^2\}$, which generate
    rotations and dilations of the symplectic plane. This is the
    semisimple sector that prevents $\mathfrak m^3/\mathfrak m^{N+1}$
    from being nilpotent.
  \item\label{lost:constants}
    The constant Hamiltonian $1\notin\mathfrak m$, on which the
    \(\hbar_{\mathrm{W}}\)-normalization of the Moyal product depends. The
    central element $\C\cdot 1\subset\mathfrak h_{\mathrm{poly}}$
    is omitted from $\mathfrak m^3$ a fortiori, so the
    \(\hbar_{\mathrm{W}}\)-normalization in the truncated theorem is conventional.
  \item\label{lost:u1-anomaly}
    The U(1) center anomaly cohomology class
    $[\bar c]\in H^2_{\mathrm{Lie}}(\bar A,\C)$ of
    \ref{thm:main-local}(f), since the cocycle
    $$\omega(z_1^kz_2^l,z_1^mz_2^n)=(kn-lm)\,
       \mathbf 1[(k+m,l+n)=(1,1)]$$
    is concentrated on multidegree $(1,1)$, with both arguments of
    total $\mathfrak m$-adic degree $1$, hence the cocycle pairs
    only elements outside $\mathfrak h^{\mathrm{nilp}}_{[3,N]}=\mathfrak m^3/
    \mathfrak m^{N+1}$. The central anomaly has support outside the
    truncation.
  \item\label{lost:colim}
    The cutoff inverse limit on the full
    $\mathfrak h$. The inverse system
    $\{\mathfrak h^{\mathrm{nilp}}_{[3,N]}\}_{N\geq 3}$ has structure maps the
    natural surjections
    $\mathfrak h^{\mathrm{nilp}}_{[3,N+1]}\twoheadrightarrow\mathfrak h^{\mathrm{nilp}}_{[3,N]}$
    induced by the $\mathfrak m$-adic projection
    \(\mathfrak m^3/\mathfrak m^{N+2}\to
    \mathfrak m^3/\mathfrak m^{N+1}\), using
    \(\mathfrak m^{N+2}\subset\mathfrak m^{N+1}\). The inverse limit
    is $\varprojlim_N\mathfrak h^{\mathrm{nilp}}_{[3,N]}=\mathfrak m^3$, the
    $\mathfrak m^3$-truncation of the Hamiltonian Lie algebra.
\end{enumerate}
What remains classically: the third-order and higher canonical
transformations, which form a nilpotent Lie subalgebra
$\mathfrak m^3\subset\mathfrak h$, a Lie ideal in $\mathfrak m^2$,
admitting the classical finite cotangent CE/PV identification
of Theorem~\ref{thm:phi-trunc-classical}.
The unreduced open BV factorization-centre morphism is governed by
separate data beyond this truncation: the cotangent centrality and
factorization-centre primitives in the unreduced open BV complex.  The
formal Moyal descendant lift is controlled by
Definition~\ref{defn:nilpotent-moyal-flat-datum}; Theorem~\ref{thm:phi-hbar-trunc}
excludes it on \(\mathfrak h^{\mathrm{nilp}}_{[3,N]}\) exactly when
\(\operatorname{cl}_{N,\hbar}\), \(\operatorname{Jac}_{N,\hbar}\), or
\(\operatorname{CE}_{N,\hbar}\) is nonzero.
\end{prop}

\begin{proof}
The quotient defining \(\mathfrak h^{\mathrm{nilp}}_{[3,N]}\) is
\(\mathfrak m^3/\mathfrak m^{N+1}\).  Therefore every monomial of total
\(\mathfrak m\)-adic degree \(1\) or \(2\), and the constant monomial,
has zero image in the quotient; this proves (i)--(iii).  The cocycle
\(\omega\) is supported on pairs of total degree \(1\), because
\(\omega(z_1^kz_2^l,z_1^mz_2^n)\) is nonzero only when
\((k+m,l+n)=(1,1)\).  Hence both inputs lie outside
\(\mathfrak m^3\), proving (iv).  For (v), a compatible family
\(\{f_N\in\mathfrak m^3/\mathfrak m^{N+1}\}_N\) gives, for each
homogeneous degree \(d\geq3\), a coefficient independent of every
window \(N\geq d\); these coefficients assemble to a unique formal
series in \(\mathfrak m^3\subset\C[[z_1,z_2]]\).  Conversely, every
element of \(\mathfrak m^3\) determines such a compatible family by
projection.  The Poisson bracket preserves the inverse-limit Lie
structure because
\(\{\mathfrak m^a,\mathfrak m^b\}\subset\mathfrak m^{a+b-2}\), so the
bracket of two compatible families is compatible in every quotient.
\end{proof}

\subsection*{Cubic inverse-limit survivor}

\begin{prop}[Cubic coordinate-window survivor and variance boundary]
\label{prop:colim-recovery}
The finite Lie quotients
\[
  \mathfrak h^{\mathrm{nilp}}_{[3,N]}=\mathfrak m^3/\mathfrak m^{N+1}
\]
recover the complete cubic Lie algebra
\(\varprojlim_N\mathfrak h^{\mathrm{nilp}}_{[3,N]}=\mathfrak m^3\).  The finite
cotangent maps \(\Phi_{\leq N}\) of
\ref{thm:phi-trunc-classical} form a compatible family of coordinate
tests.  Their functoriality is contravariant: it is pullback along Lie
quotients, with additional data required for an inverse system of dg CE
cochain projections.  A Lie quotient
\(\mathfrak h^{\mathrm{nilp}}_{[3,M]}\twoheadrightarrow\mathfrak h^{\mathrm{nilp}}_{[3,N]}\) induces CE
cochain maps by pullback in the opposite direction.  Consequently
\(\Phi_{\mathfrak m^3}\) denotes the coordinate-window family
\(\{\Phi_{\leq N}\}_N\) together with the chain-side inverse limit and an
exact continuous-dualization datum; the expression
\(\varprojlim_N\Phi_{\leq N}\) on CE cochains denotes that datum.

This is the maximal survivor of the nilpotent tower only in that precise
sense: the Lie-side tower recovers \(\mathfrak m^3\), while a completed
cochain \(P_0\)-statement is governed by the variance and
bracket-admissibility data named above.
\end{prop}

\begin{proof}
The Lie-side inverse limit is
\[
  \varprojlim_N\mathfrak m^3/\mathfrak m^{N+1}=\mathfrak m^3
\]
with its \(\mathfrak m\)-adic topology: the coefficient of each
monomial of total degree \(d\ge3\) stabilizes in every quotient with
\(N\ge d\), and these stabilized coefficients are exactly a formal
series with no terms of degree \(<3\).  The variance obstruction appears already at
\(M=4\), \(N=3\).  In \(\mathfrak h^{\mathrm{nilp}}_{[3,4]}\),
\[
  \{z_1^3,z_2^3\}=9z_1^2z_2^2,
  \qquad
  d_{\mathrm{CE}}c^{2,2}\supset -9c^{3,0}c^{0,3}.
  \]
The coordinate projection to the \(N=3\) cochain window kills
\(c^{2,2}\) and leaves \(c^{3,0}c^{0,3}\), so
\[
  p(d_{\mathrm{CE}}c^{2,2})=-9c^{3,0}c^{0,3}
  \neq0=d_{\mathrm{CE}}p(c^{2,2}).
  \]
Thus projection of CE cochains has a nonzero dg defect.  The correct finite
functoriality is contravariant pullback along Lie quotients, while a
completed cochain theorem must be obtained either from a chain-side
inverse limit followed by exact continuous dualization or from a
separately given topology whose cochain transition maps are dg maps.

The Lie algebra $\mathfrak m^3$ is lower-central Tate-pronilpotent in the
$\mathfrak m$-adic Tate filtration: its descending Lie powers satisfy
$(\mathfrak m^3)^{(\geq m)}\subset\mathfrak m^{m+3}$ because each
iterated bracket raises the $\mathfrak m$-adic degree by at least
$1$ when both arguments lie in $\mathfrak m^3$ (since
$\{H_p,H_q\}\subset\mathfrak m^{p+q-2}$ and $p,q\geq 3$ implies
$p+q-2\geq p+1$). Hence $\bigcap_m(\mathfrak m^3)^{(\geq m)}=0$,
	and the lower-central Tate-pronilpotence hypothesis of
\ref{lem:continuous-bar-cobar} holds. The Tate Casimir
$C_{\mathfrak m^3}=\sum_{d\geq 3}\sum_i e_{d,i}\otimes e_d^i$
is governed by the same product/direct-sum convergence question
as in \ref{lem:tate-casimir-obstruction}, but on the smaller
	direct-sum dual $\bigoplus_{d\geq 3}H_d^\vee\subset
\bigoplus_{d\geq 1}H_d^\vee$.  The finite maps
\(\Phi_{\leq N}\) at every finite \(N\) use only the finite Casimir
\(\sum_{3\leq d\leq N}\sum_i e_{d,i}\otimes e_d^i\).  Passing from these
finite identities to a completed cochain identity is exactly the
continuous-dualization and variance datum isolated in the statement, rather
than a formal inverse-limit consequence.
\end{proof}

\begin{defn}[Low-degree obstruction datum for the nilpotent tower]
\label{defn:nilpotent-low-degree-obstruction-terms}
  For \(N\geq3\) let
  \[
    \pi_N\colon\mathfrak m^3\twoheadrightarrow
    \mathfrak h^{\mathrm{nilp}}_{[3,N]}=\mathfrak m^3/\mathfrak m^{N+1}
  \]
  be the nilpotent window projection.  The two obstruction terms for
  adjoining the removed low-degree sector are:
  \[
    \mathfrak o^{\mathrm{lin}}_N(x;\tau)
      =\{x,\tau\}\bmod \mathfrak m^{N+1},
      \qquad
      x\in H_1,\ \tau\in \mathfrak m^{N+1}/\mathfrak m^{N+2},
  \]
  and
  \[
    \mathfrak o^{\mathrm{ss}}\colon\Lambda^2H_2\to H_2,\qquad
    \mathfrak o^{\mathrm{ss}}(x,y)=\{x,y\}.
  \]
  Concretely
  \[
    \mathfrak o^{\mathrm{lin}}_N(z_1;z_2^{N+1})
      =(N+1)z_2^N\neq0
      \quad\text{in }\mathfrak m^N/\mathfrak m^{N+1},
  \]
  and
  \[
    \{z_1^2,z_2^2\}=4z_1z_2,\qquad
    \{z_1^2,z_1z_2\}=2z_1^2,\qquad
    \{z_1z_2,z_2^2\}=2z_2^2 .
  \]
  The compatible Casimir missing from the nilpotent tower is
  \[
    \mathfrak o^{\mathrm{Cas}}_{\leq2}
      =\sum_i e_{1,i}\otimes e_1^i
       +\sum_i e_{2,i}\otimes e_2^i .
  \]
\end{defn}

\begin{lem}[Obstruction-datum criterion for enlarging the tower]
\label{lem:nilpotent-obstructions-exact}
  A finite \(\mathfrak m\)-adic tower extending
  \(\{\mathfrak h^{\mathrm{nilp}}_{[3,N]}\}_{N\geq3}\) to a Lie tower containing
  \(H_1\oplus H_2\) and using the same quotient maps exists only if
  \(\mathfrak o^{\mathrm{lin}}_N=0\) for all \(N\) and
  \(\mathfrak o^{\mathrm{ss}}\) is nilpotent in the lower-central
  filtration.  A cotangent CE/PV extension of the inverse limit to the
  full \(\mathfrak h=\mathfrak m\) also requires a coefficient category
  in which the compatible tensor
  \(\mathfrak o^{\mathrm{Cas}}_{\leq2}
    +\sum_{d\geq3}\sum_i e_{d,i}\otimes e_d^i\)
  converges.
\end{lem}

\begin{proof}
  If \(H_1\) is included and the window \(N\) is still the quotient by
  \(\mathfrak m^{N+1}\), then the tail
  \(\mathfrak m^{N+1}\) must be a Lie ideal for the enlarged source.
  The displayed value
  \(\{z_1,z_2^{N+1}\}=(N+1)z_2^N\notin\mathfrak m^{N+1}\)
  shows that this ideal condition is exactly the vanishing of
  \(\mathfrak o^{\mathrm{lin}}_N\), and the displayed value is nonzero.

  If \(H_2\) is included, then its bracket remains inside \(H_2\) and
  equals the nonabelian \(\mathfrak{sl}_2\) bracket displayed in
  Definition~\ref{defn:nilpotent-low-degree-obstruction-terms}.  The
  image of \(\mathfrak o^{\mathrm{ss}}\) is all of \(H_2\), so the
  descending Lie powers never leave \(H_2\).  This is exactly the
  obstruction to nilpotence and to the conilpotent bar-cobar hypothesis.

  The finite cotangent windows for
  \(\mathfrak h^{\mathrm{nilp}}_{[3,N]}\) use the finite Casimir
  \(\sum_{3\leq d\leq N}\sum_i e_{d,i}\otimes e_d^i\).  Extending to
  \(\mathfrak h\) adds the missing \(d=1,2\) components and then the
  whole compatible family over all \(d\geq1\).  Thus the full
  cotangent lift requires precisely the convergence of the displayed
  full Casimir family.
\end{proof}

\begin{thm}[Low-degree obstruction to nilpotent-window recovery of the full Hamiltonian source]
\label{thm:colim-obstruction-full-h}
The coordinate-window family
  $\Phi_{\mathfrak m^3}$ of \ref{prop:colim-recovery} is the
  classical cotangent lift datum on $\mathfrak m^3$. Recovery of all of
$\mathfrak h=\mathfrak m$ must
include the low-degree sectors
$\mathfrak m/\mathfrak m^3=H_1\oplus H_2=\C\cdot z_1\oplus
\C\cdot z_2\oplus\mathrm{span}\{z_1^2,z_1z_2,z_2^2\}$, which by
\ref{lem:finite-window-hamiltonian-obstruction}(iii) and
\ref{lem:nilpotent-obstructions-exact} has a stability obstruction under
any $\mathfrak m$-adic cutoff: the bracket
$\{z_1,z_2^{N+1}\}=(N+1)z_2^N$ is the explicit obstruction.
Equivalently, the obstruction datum is
\[
  \bigl(\mathfrak o^{\mathrm{lin}}_N,\,
    \mathfrak o^{\mathrm{ss}},\,
    C_{\mathfrak h}\bigr),
\]
where
\[
  \mathfrak o^{\mathrm{lin}}_N(z_1;z_2^{N+1})
    =(N+1)z_2^N\neq0,\qquad
  \operatorname{im}\mathfrak o^{\mathrm{ss}}=H_2\cong\mathfrak{sl}_2,
\]
and \(C_{\mathfrak h}\) is the full infinite Casimir family below.
Equivalently, the Tate Casimir
$C_{\mathfrak h}=\sum_{d\geq 1}\sum_i e_{d,i}\otimes e_d^i$ has
low-degree components
$\sum_i e_{1,i}\otimes e_1^i\in H_1\otimes H_1^\vee$ and
$\sum_i e_{2,i}\otimes e_2^i\in H_2\otimes H_2^\vee$
that the truncation removes; the degree-$1$ component is the source
of the linear-Hamiltonian translation, and the degree-$2$ component
is the source of the conformal-symplectic $\mathfrak{sl}_2$ direction,
where semisimplicity replaces nilpotence. Recovering the full
$\mathfrak h$ requires a coefficient category in which
$C_{\mathfrak h}$ converges \emph{including} both low-degree
components, which is exactly the Casimir support obstruction of
\ref{lem:tate-casimir-obstruction}. The truncation route leaves that
obstruction intact; it removes the sector supporting the
obstruction.  A recovery theorem using this nilpotent-window route must
therefore give
both a coefficient category in which \(C_{\mathfrak h}\) converges and a
relative or reductive model for the \(H_2\)-sector.
\end{thm}

\begin{proof}
  The tower used in \ref{prop:colim-recovery} is
  \(\mathfrak h^{\mathrm{nilp}}_{[3,N]}=\mathfrak m^3/\mathfrak m^{N+1}\).
  The transition map
  \(\pi_{N+1,N}\colon\mathfrak m^3/\mathfrak m^{N+2}\to
  \mathfrak m^3/\mathfrak m^{N+1}\) is induced by the inclusion
  \(\mathfrak m^{N+2}\subset\mathfrak m^{N+1}\).  Therefore
  \[
    \varprojlim_N \mathfrak h^{\mathrm{nilp}}_{[3,N]}
      =\varprojlim_N \mathfrak m^3/\mathfrak m^{N+1}
      =\mathfrak m^3
  \]
  as a complete filtered Lie algebra.  The homogeneous pieces
  \(H_1\) and \(H_2\) have zero image in every term of this tower;
  hence the inverse limit misses them.

  The obstruction is intrinsic to the chosen finite windows.  Including
  \(H_1\) while keeping the same \(\mathfrak m\)-adic finite windows gives
  a Lie-ideal obstruction in the tails: for every \(N\),
  \[
    z_2^{N+1}\in\mathfrak m^{N+1},\qquad
    \{z_1,z_2^{N+1}\}=(N+1)z_2^N\notin\mathfrak m^{N+1}.
  \]
  Therefore the quotient by the tail carries a full-Hamiltonian bracket
  obstruction.  If one includes \(H_2\), the degree-two
  Hamiltonians contain the \(\mathfrak{sl}_2\) bracket
  \([H_2,H_2]=H_2\), so the finite quotient violates nilpotence.  Thus a
  finite nilpotent \(\mathfrak m\)-adic tower is incompatible with the
  low-degree sector \(H_1\oplus H_2\).

  Finally, the full Casimir \(C_{\mathfrak h}\) has nonzero components
  in \(H_1\otimes H_1^\vee\), \(H_2\otimes H_2^\vee\), and in every
  higher \(H_d\otimes H_d^\vee\).  The nilpotent tower sees only the
	  \(d\geq3\) part.  Hence the inverse-limit cotangent lift on
	  \(\mathfrak m^3\) reaches the full Hamiltonian sector only after
	  an additional coefficient category makes the full Casimir converge.
\end{proof}

\begin{cor}[Full maximal-ideal obstruction class]
\label{cor:full-maximal-ideal-universal-obstruction-class}
  For the full reduced formal Hamiltonian algebra
  \(\mathfrak h=\mathfrak m=(z_1,z_2)\C[[z_1,z_2]]\), the direct
  formal-local CE/PV recognition theorem remains valid in the
  finite-coordinate product completion.  The filtered-cobar
  Koszul recognition criterion of
  Theorem~\ref{thm:universal-filtered-cobar-recognition}, however, has
  a nonzero obstruction class in
  \(\mathcal O_{\mathrm{Kos}}\) of
  Definition~\ref{defn:universal-koszul-obstruction-complex}.  Its
  visible components are
  \[
    \mathfrak o^{\mathrm{lin}}_N(z_1;z_2^{N+1})
      =(N+1)z_2^N\neq0,
    \qquad
    \operatorname{im}\mathfrak o^{\mathrm{ss}}=H_2\cong\mathfrak{sl}_2,
  \]
  \[
    [C_{\mathfrak h}]
      =
      \left[
        \sum_{d\geq1}\sum_i e_{d,i}\otimes e_d^i
      \right]
      \in Q_{\mathrm{Cas}},
    \qquad
    H_2\subset\bigcap_{r\geq1}\mathfrak h^{(r)} .
  \]
  Therefore a full-\(\mathfrak h\)
  bar-cobar/enveloping theorem requires,
  before any open-side \(A_\infty\) comparison, all of the following:
  a coefficient topology representing \(C_{\mathfrak h}\), continuity of
  the finite-window bracket and coadjoint transition maps, vanishing
  \(\varprojlim^1\) for the completed CE/PV and cobar towers, a
  relative or reductive treatment of the
  \(H_2\)-sector, and a twisting cochain whose associated-graded cobar
  map has acyclic cone.
\end{cor}

\begin{proof}
  The formal-Darboux CE/PV part is
  Theorem~\ref{thm:formal-disk-completed-ce-pv}, equivalently the
  formal-Darboux instance of
  Theorem~\ref{thm:universal-formal-local-ce-pv-recognition}: the
  bracket and coadjoint formulas are continuous for finite coordinate
  projections, so the generator substitution
  \(c^\rho\mapsto\partial_\rho\), \(u_f\mapsto O_f\), and
  \(d_{\mathrm{CE}}u_f\mapsto X_f=d_\pi O_f\) is a completed dg
  commutative cochain isomorphism after the \(u\)-regrading.  The
  standard Chevalley--Eilenberg complex retains \(|u_f|_{\mathrm{CE}}=2\);
  the coordinate complex
  \(\operatorname{Regr}_uC^\bullet_{\mathrm{CE},\mathrm{coord}}\)
  has \(|u_f|_{\mathrm{coord}}=0\).  The \(P_0\)-algebra reading is finite-level or
  bracket-admissible data; an unqualified inverse limit of ambient
  coordinate products gives a different object.

  The filtered-cobar recognition theorem has an additional linear
  obstruction, the same computation as in
  Definition~\ref{defn:nilpotent-low-degree-obstruction-terms}:
  adjoining \(H_1\) to the \(\mathfrak m\)-adic nilpotent windows makes
  the tail non-invariant because
  \(\{z_1,z_2^{N+1}\}=(N+1)z_2^N\notin\mathfrak m^{N+1}\).  The
  semisimple obstruction is the displayed \(\mathfrak{sl}_2\) bracket
  on \(H_2\), hence \(H_2\subset\bigcap_r\mathfrak h^{(r)}\).  This
  gives a nonzero lower-central component in
  \(\mathcal O_{\mathrm{conilp}}\).

  The Casimir component is nonzero in the strict product/direct-sum
  topology by the support argument of
  Lemma~\ref{lem:tate-casimir-obstruction}: the compatible
  finite-window Casimirs have nonzero identity component in every
  \(H_d\otimes H_d^\vee\), while a strict tensor with a direct-sum dual
  has only finitely many dual-degree components.  Finally, even after
  changing topology and handling \(H_2\) relatively, the
  filtered-cobar criterion still demands the cobar cone and the Milnor
  \(\varprojlim^1\) terms in
  Definition~\ref{defn:universal-koszul-obstruction-complex} to vanish.
  These are precisely the listed remaining requirements.
\end{proof}

\begin{prob}[Relative replacement of the nilpotent
\(\mathfrak m^3\)-window]
\label{prob:linear-sector-lift}
Replace the coordinate-window Koszul-resolution datum
$\Phi_{\mathfrak m^3}$ of \ref{prop:colim-recovery} from
$\mathfrak m^3$ by a relative construction on all of
$\mathfrak h=\mathfrak m$ by giving:
(D1) a Tate or weighted-coefficient extension in which the degree-$1$
and degree-$2$ Casimir components converge as kernels for the BV
propagator; (D2) a relative or reductive replacement for the
lower-central Tate-pronilpotence hypothesis of
\ref{lem:continuous-bar-cobar} in which the conformal-symplectic
$\mathfrak{sl}_2$ direction $H_2$ is included as semisimple data.  By
\ref{defn:nilpotent-low-degree-obstruction-terms} the obstruction
terms are the triple
\(\bigl(\mathfrak o^{\mathrm{lin}}_N,
   \mathfrak o^{\mathrm{ss}},
   \mathfrak o^{\mathrm{Cas}}_{\leq2}\bigr)\);
by \ref{lem:tate-casimir-obstruction} (D1) is the same
coefficient-category problem as
\ref{prob:formal-factorization-center}.  The nilpotent-truncation route
closes the Koszul resolution on the cubic ideal \(\mathfrak m^3\);
the displayed triple is the precise obstruction to using the same
pronilpotent route for the full \(\mathfrak h\).  It is not an
  obstruction to the completed formal-Darboux CE/PV construction of
Theorem~\ref{thm:main-local}.
\end{prob}

\begin{thm}[Semisimple obstruction to lower-central Tate-pronilpotence]
\label{thm:conilpotency-essential}
A second obstruction to extending the nilpotent \(\mathfrak m^3\)
bar-cobar route to all of \(\mathfrak h\) is structural: the inclusion of
$H_2=\mathfrak{sl}_2$ violates the lower-central Tate-pronilpotence route
used in Lemma~\ref{lem:continuous-bar-cobar} because
$[H_2,H_2]=H_2$ is its own descending Lie series. Even if the
degree-$1$ and degree-$2$ Casimir components converged in some
weighted Tate completion, direct application of the pronilpotent
bar-cobar route would still be obstructed.  Ordinary finite-dimensional
bar-cobar and PBW formalism remain valid for \(\mathfrak{sl}_2\); the
lower-central Tate-pronilpotent criterion must be replaced, on a
\(L\) containing \(H_2\), by a relative or reductive model that builds in
semisimple subalgebras separately.
\end{thm}

\begin{proof}
  Let \(L\) be any filtered Lie algebra containing the degree-two
  Hamiltonians \(H_2\) as a Lie subalgebra with the usual Poisson
  bracket.  The descending Lie powers of \(L\) satisfy
  \(L^{(n)}\cap H_2\supset H_2\) for every \(n\), because
  \([H_2,H_2]=H_2\).  Hence
  \[
    H_2\subset \bigcap_{n\geq1} L^{(n)} .
  \]
  The intersection of the descending Lie powers is therefore nonzero.
  This contradicts the lower-central Tate-pronilpotence condition used in
  Lemma~\ref{lem:continuous-bar-cobar}.  The bar-cobar argument applies
  to the nilpotent radical \(\mathfrak m^3\).  A Lie source containing
  the semisimple \(H_2\)-sector requires a relative or reductive
  replacement.
\end{proof}

\subsection*{Costello--Li precedent}

\begin{rmk}[Costello--Li balanced matrix truncation as structural
model]\label{rmk:costello-li-precedent}
The Costello--Li BCOV construction
\cite{costello-li-quantum-bcov}*{\S3} supplies the
holomorphic-volume precedent for a
superrank-stable connected trace construction through the open gauge
algebra
$\mathfrak{gl}(N\mid N)$, with the directed Lie embedding
$\mathfrak{gl}(N\mid N)\hookrightarrow\mathfrak{gl}(N+k\mid N+k)$
and the BV-cohomology vanishing
$H^*_{\mathrm{BV}}\bigl(\mathfrak{gl}(N\mid N)\bigr)=0$
(supertrace identity $\mathrm{str}(I)=N-N=0$ kills the putative
U(1)-anomaly). The structural analogy with the truncation here:
\begin{enumerate}
  \item Each window $\mathfrak{gl}(N\mid N)$ is finite-dimensional,
    BV-cohomology vanishes, the universal CE/Schouten identification
    and the BV-propagator coefficient kernel are well-defined on
    each window.
  \item The inverse / direct system is the Lie embedding
    $\mathfrak{gl}(N\mid N)\hookrightarrow\mathfrak{gl}(N+1\mid N+1)$
    on the Costello--Li side, the Lie surjection
    $\mathfrak h^{\mathrm{nilp}}_{[3,N+1]}\twoheadrightarrow\mathfrak h^{\mathrm{nilp}}_{[3,N]}$
    here.
  \item The U(1)-anomaly is killed by the structure of each window
    on the Costello--Li side (supertrace), and is removed by the
    $\mathfrak m^3$-truncation here (the cocycle $\omega$
    pairs only multidegree-$(1,0)$ with multidegree-$(0,1)$, both
    in $H_1\subset\mathfrak m\setminus\mathfrak m^2$, hence outside
    $\mathfrak m^3$).
  \item The superrank colimit converges on the Costello--Li side
    to the full $\mathfrak{gl}(\infty\mid\infty)$, providing the
    superrank-stable connected single-trace algebra. Here the inverse limit
    converges to the proper subalgebra
    $\mathfrak m^3\subsetneq\mathfrak h$ by
    Theorem~\ref{thm:colim-obstruction-full-h}.
\end{enumerate}
The structural difference: the
$\mathfrak{gl}(N\mid N)$ system is graded by matrix size, with
trivial action of one window on the next-window orthogonal
complement; the
$\mathfrak h=\mathfrak m$ system is graded by $\mathfrak m$-adic
degree, with the linear and quadratic sectors acting non-trivially
on every higher-degree window, which is precisely the
finite-window obstruction of
\ref{lem:finite-window-hamiltonian-obstruction}(iii) plus the
$\mathfrak{sl}_2\subset H_2$ semisimple obstruction to
lower-central Tate-pronilpotence.
\end{rmk}
